\documentclass[10pt,a4paper]{amsart}
\usepackage[margin=19mm]{geometry}
\usepackage{amsmath,amssymb,mathtools}
\usepackage[scr=rsfs,cal=euler]{mathalfa}
\usepackage{xfrac}
\usepackage[unicode,hidelinks,bookmarks=false]{hyperref}
\newcommand{\ie}{i.e.\ }
\newcommand{\cf}{cf.\ }
\newcommand{\resp}{respectively\ }
\newcommand{\Subsection}{\subsection}
\providecommand{\nobreakdash}{\nobreak\hbox{-}}
\setlength{\emergencystretch}{2em}
\multlinegap0pt
\usepackage{bm}
\usepackage{float}
\usepackage{amssymb}
\usepackage{algorithm}\usepackage{algpseudocode}
\usepackage{tikz}
\usepackage{mathtools}
\usepackage{cancel}
\usepackage{xcolor}
\usepackage[normalem]{ulem}
\newtheorem{lemma}{Lemma}
\newtheorem{theorem}{Theorem}
\newcommand{\FOMsize}{N}
\newcommand{\Jcheck}{\check{\bmJ}}
\newcommand{\R}{\mathbb{R}}
\newcommand{\ROMsize}{r}
\newcommand{\Rcheck}{\check{\bmR}}
\newcommand{\bmB}{\bm{B}}
\newcommand{\bmJ}{\bm{J}}
\newcommand{\bmK}{\bm{K}}
\newcommand{\bmR}{\bm{R}}
\newcommand{\bmU}{\bm{U}}
\newcommand{\bmW}{\bm{W}}
\newcommand{\bzero}{\mathbf{0}}
\newcommand{\matset}{\mathbf{S}}
\newcommand{\portsize}{m}
\newcommand{\stateapprox}{\tilde{\bm x}}
\newcommand{\statex}{{\bm x}}
\title{Structure-Preserving Generalized Manifold Galerkin Reduction for Port-Hamiltonian Systems}
\author{Silke Glas, Hongliang Mu}
\date{Source: arXiv:2603.08656v2 (CC BY 4.0)}
\begin{document}
\maketitle

\noindent\emph{Source and licence.} Structure-Preserving Generalized Manifold Galerkin Reduction for Port-Hamiltonian Systems. Version arXiv:2603.08656v2 (CC BY 4.0), distributed by arXiv under the Creative Commons Attribution 4.0 International licence, http://creativecommons.org/licenses/by/4.0/, as recorded in the arXiv licence metadata for this exact version and shown on the abstract page https://arxiv.org/abs/2603.08656v2. Full licence notice: \texttt{licenses/arxiv-2603.08656-CC-BY-4.0.txt}.\par\smallskip

\noindent\textit{Editorial scope.} Counted unit: the article's theorem Theo:exists\_JR\_inv\_BVbar (source0.tex lines 737--765). The article's complete original proof of it is delivered with it (source0.tex lines 766--898, one \texttt{proof} environment of 133 lines). No mathematics is written, repaired or re-derived: the statement and the proof are the article's own lines, in the article's own notation, and no retained line is substituted or re-flowed.  Retained uncounted as the source-local context the proof needs: lines 675--678; lines 683--687; lines 720--723.  Nothing was omitted from any retained span, so the package carries no omission record.  No retained line contains a citation, so the wrapper carries no bibliography and no bibliography entry.  No retained line contains a figure environment, an image-inclusion command or drawing code, so no image is delivered and the package declares no asset dependency.  Editorial formatting only, in addition: an AMS \texttt{amsart} preamble that restates the article's own declarations and macros used by the retained lines, namely the environments \texttt{lemma}, \texttt{theorem} on separate counters, together with the standard packages those lines need.  The retained environments print in this document as Lemma 1, Theorem 1 in order of appearance, whereas the article prints the same items with the numbering of its own full document.  The complete native source of the article as distributed in the arXiv e-print for this exact version is bundled unchanged as \texttt{sources/196021/source0.tex}.
 \begin{lemma}
		\label{lemma:analytic_measure}
		Let $f:\R^{\FOMsize}\rightarrow\R$ be a real analytic function and let $m_{\FOMsize}$ be the Lebesgue measure on $\R^{\FOMsize}$. If $m_{\FOMsize}(\{x\in\R^{\FOMsize}| f(x)=0\})>0$, then $f\equiv0$.
	\end{lemma}

\begin{equation}
		\label{equ:GMG_defi_Jcheck_Rcheck}
				\Jcheck \coloneq \frac{1}{2}\left((\bmJ-\bmR)^{-1} - (\bmJ-\bmR)^{-\top}\right)\ \quad \mathrm{and}\quad\Rcheck \coloneq
				\frac{1}{2} \left(-(\bmJ-\bmR)^{-1} -(\bmJ-\bmR)^{-\top}\right).
			\end{equation}

\begin{theorem}
\label{Theo:exists_JR_inv}
			Let $\bmJ=-\bmJ^{\top}\in\R^{\FOMsize\times \FOMsize}$ be a skew-symmetric matrix and let $\bmR=\bmR^{\top}\in\R^{\FOMsize\times\FOMsize}$ be a positive semi-definite matrix. If $\bmR\neq\bzero$, i.e., $\bmR$ is not the zero matrix, and the matrix $(\bmJ-\bmR)$ is invertible, then for almost every matrix $\bmU\in \R^{\FOMsize\times\ROMsize}, 1\leq\ROMsize\leq\FOMsize$, it holds that $\bmU\in\matset_{(\bmJ-\bmR)^{-1}}$, i.e., $\mathrm{det}\left(\bmU^{\top}(\bmJ-\bmR)^{-1}\bmU\right)\neq 0$.
\end{theorem}	

\begin{theorem}
\label{Theo:exists_JR_inv_BVbar}
Let $\bmJ=-\bmJ^{\top}\in\R^{\FOMsize\times \FOMsize}$ be skew-symmetric, let $\bmR=\bmR^{\top} \in\R^{\FOMsize\times\FOMsize}$ be positive semi-definite with $\bmR \neq \bm{0}$ and assume that $(\bmJ-\bmR)$ is invertible. Let $\bmB\in\R^{\FOMsize\times\portsize}$ be of full column rank and define 
\begin{align*}
	\ell := \mathrm{rank}\left(\bmB, (\bmJ-\bmR)^{-1}\bmB \right).
\end{align*}
Furthermore, let $\Jcheck,\Rcheck \in \R^{\FOMsize \times \FOMsize}$ be defined as in \eqref{equ:GMG_defi_Jcheck_Rcheck}. Assume that 
\begin{align} \label{eq:full_space_assumption}
			 \mathrm{ker}(\Rcheck)\cup
			 \mathrm{span}(\bmB)\cup
			 \mathrm{span}((\bmJ-\bmR)^{-1}\bmB)
			 \neq
			 \R^{\FOMsize},
\end{align}			 
and that the matrices $\bmB^{\top}(\bmJ-\bmR)^{-1}\bmB$ and $\bmB^{\top}\bmR\bmB$ are invertible. Then, for almost every matrix $\bmU\in \R^{\FOMsize\times\ROMsize} (\ROMsize\leq\FOMsize-\ell)$, it holds that 
\begin{align*}
	\begin{bmatrix}
				\bmB,\bmU
			\end{bmatrix}\in\matset_{(\bmJ-\bmR)^{-1}},
\end{align*}
i.e., 
\begin{align*}
\mathrm{det}\left(\begin{bmatrix}
				\bmB,\bmU
			\end{bmatrix}^{\top}(\bmJ-\bmR)^{-1}\begin{bmatrix}
				\bmB,\bmU
			\end{bmatrix}\right)\neq 0.
\end{align*}
\end{theorem}

\begin{proof}
For any $1\le\ROMsize\leq \FOMsize-\ell$, let $f_{\mathrm{det},\ROMsize,2}\colon\R^{\FOMsize\times\ROMsize}\rightarrow \R$ be defined by 
			\begin{align*}
				f_{\mathrm{det},\ROMsize,2}(\bmU)\coloneq \mathrm{det}\left(\begin{bmatrix} \bmB,\bmU \end{bmatrix}^{\top} \left(\Jcheck-\Rcheck\right) \begin{bmatrix} \bmB,\bmU \end{bmatrix}\right).
			\end{align*}
We note that $f_{\mathrm{det},\ROMsize,2}$ is a polynomial with respect to the entries of $\bmU$, i.e. $f_{\mathrm{det},\ROMsize,2}$ is a real analytic function. By Lemma~\ref{lemma:analytic_measure}, it is sufficient to show that $f_{\mathrm{det},\ROMsize,2}$ is not identically zero. Equivalently, it remains to construct a matrix  $ \bmU\in\R^{\FOMsize\times \ROMsize} $ for $\ROMsize=1,\ldots,\FOMsize-\ell$ such that
			\begin{equation*}
				\mathrm{det}\left(\begin{bmatrix} \bmB,\bmU \end{bmatrix}^{\top} \left(\Jcheck-\Rcheck\right) \begin{bmatrix} \bmB,\bmU \end{bmatrix}\right) \neq 0. 
\end{equation*}
Let $\bmW\in\R^{\FOMsize\times(\FOMsize-\ell)}$ be a matrix whose columns form an orthogonal complement of \\$\textnormal{span}(\bmB,(\Jcheck-\Rcheck)\bmB)$. 
			This means we have 
			\begin{align}
				\label{equ:exists_JR_inv_BVbar_orth_comp}
				\bmB^\top\bmW = \bm{0},\
				\bmB^\top (\Jcheck-\Rcheck)^\top \bmW = \bm{0},\ \textnormal{ and }
				\mathrm{span}(
					\bmB, (\Jcheck-\Rcheck)\bmB, \bmW
				) =\R^{\FOMsize}.
			\end{align}
			Define
			$
			\bmJ_1 \coloneq  \bmW^{\top}\Jcheck\bmW$ and
			$\bmR_1 \coloneq  \bmW^{\top}\Rcheck\bmW.
			$
			Further, assume that $\mathrm{det}(\bmJ_1-\bmR_1)= 0$. Then, there exists a nonzero vector $\statex\in\R^k$ s.t.
				$(\bmJ_1-\bmR_1)\statex = \bzero$.
			From the latter equation, we obtain
			\begin{equation}
				\label{equ:GMG_exist_Orth_cond}
				\bmW^\top (\Jcheck-\Rcheck)\bmW\statex =\bm{0}\ \textnormal{ and }
				\statex^{\top}\bmW^{\top}\Rcheck\bmW\statex\allowbreak= 0.
			\end{equation}
			From \eqref{equ:exists_JR_inv_BVbar_orth_comp} and \eqref{equ:GMG_exist_Orth_cond}, we get 
			\begin{align} \label{eq:orthcond}(\Jcheck-\Rcheck)^{\top}\bmW\perp \textnormal{span}\left( \bmW\statex, \bmB, (\bmJ-\bmR) \bmB\right), \end{align}
			where we used used
			\begin{align*}
				\bm{0} = \bmB^\top\bmW = \bmB^\top (\bmJ-\bmR)^\top (\bmJ-\bmR)^{-\top}  \bmW = \bmB^\top (\bmJ-\bmR)^\top (\Jcheck-\Rcheck)^{\top}  \bmW.
			\end{align*}
			We have that the dimension of the space $\textnormal{span}(\bmB,(\Jcheck-\Rcheck)\bmB) = \textnormal{span}(\bmB,(\bmJ-\bmR)^{-1}\bmB)$. Then, due to $\textnormal{det}(\bmJ-\bmR)\neq 0$, it follows that the dimension of the space $\textnormal{span}((\bmJ-\bmR)\bmB, \bmB)$ is $\ell$. Further, since the dimension of the space $\textnormal{span} ((\Jcheck-\Rcheck)^{\top}\bmW)$ is $\FOMsize-\ell$, it follows due to the orthogonality conditions \eqref{eq:orthcond} that 				
			$$
				\mathrm{rank}\left(
				\begin{bmatrix}
					(\bmJ-\bmR)\bmB, \bmB, (\Jcheck-\Rcheck)^{\top}\bmW
				\end{bmatrix}
				\right) =\FOMsize.$$
				From the latter rank argument combined with \eqref{eq:orthcond}, we obtain $\bmW\statex\in\mathrm{span}\left((\bmJ-\bmR)\bmB\right)$. 
			
			If $\ell=\portsize$, then
			  we get $
			\mathrm{dim}\left(\textnormal{span}
				((\bmJ-\bmR)\bmB, \bmB))\right)
			=\portsize.
			$ 
			Since the matrix $\bmB$ is assumed to have full column rank, we arrive at $\bmW\statex\in \mathrm{span}
			\left((\bmJ-\bmR)\bmB\right)\subseteq
			\mathrm{span}(\bmB).
			$
			This yields $\bmW\statex\in \textnormal{span}(\bmW)\cap \textnormal{span}( \bmB)= \{ \bm{0} \}.$
			Thus, we have $\bmW\statex=\bm{0}$, yielding $\statex=\bm{0}$. However, this is a contradiction since we assume $\statex\neq\bm{0}$ and it follows that $\mathrm{det}(\bmJ_1-\bmR_1)\neq0. $
			
		    We now consider the remaining case $\ell\neq\portsize$.
		    We know that  $\exists\stateapprox\in\R^{\portsize}$ s.t. $\bmW\statex=(\bmJ-\bmR)\bmB\stateapprox$ and with \eqref{equ:GMG_exist_Orth_cond} we arrive at
			\begin{align*}
				0 =&\statex^{\top}\bmW^{\top}\Rcheck\bmW\statex
				=\stateapprox^{\top}\bmB^{\top}(\bmJ-\bmR)^{\top}
				\Rcheck
				(\bmJ-\bmR)\bmB\stateapprox\\
				=&
				-\stateapprox^{\top}\bmB^{\top}(\bmJ-\bmR)^{\top}
				(\Jcheck-\Rcheck)
				(\bmJ-\bmR)\bmB\stateapprox\\
				=&
				-\stateapprox^{\top}\bmB^{\top}(\bmJ-\bmR)^{\top}
				\bmB\stateapprox
				=
				\stateapprox^{\top}\bmB^{\top}\bmR
				\bmB\stateapprox.
			\end{align*}
			Since $\bmB^{\top}\bmR\bmB$ is non-degenerate and $\bmW$ is of full column rank, we get $\statex= \bzero$, which is again a contradiction to our assumption of a nonzero $\statex$. Also in this case we arrive at $\mathrm{det}(\bmJ_1-\bmR_1)\neq0.$
			
Recall the assumption \eqref{eq:full_space_assumption}. Since $\textnormal{span}(\bmW)$ is by construction the orthogonal complement of $\mathrm{span}(
					\bmB, (\Jcheck-\Rcheck)\bmB)$ \eqref{equ:exists_JR_inv_BVbar_orth_comp}, it follows that $\textnormal{span}(\bmW) \not\subseteq \mathrm{ker}(\Rcheck)$. Thus, there exists $\bm{w} \in \bmW$ such that $\bm{w}^{\top} \Rcheck \bm{w} >0$ and we obtain $\bmR_1 = \bmW^{\top}\Rcheck\bmW \neq \bzero$.	
					
Since $(\bmJ_1-\bmR_1)$ is non-degenerate and $\bmR_1 \neq \bzero$, we can construct $\Jcheck_1,\Rcheck_1\in\R^{(\FOMsize-\ell)\times(\FOMsize-\ell)}$ based on \eqref{equ:GMG_defi_Jcheck_Rcheck} s.t.
			\begin{equation*}
				\Jcheck_1-\Rcheck_1 = (\bmJ_1-\bmR_1)^{-1},\quad
				\Jcheck_1=-\Jcheck_1^{\top},\quad
				\Rcheck_1 =\Rcheck_1^{\top}\succeq 0,\quad
				\Rcheck_1\neq \bzero.				
			\end{equation*}
By applying Theorem \ref{Theo:exists_JR_inv} to $\Jcheck_1, \Rcheck_1$, we know that for almost every matrix $\bmK\in\R^{(\FOMsize-\ell)\times \ROMsize}$, $1 \leq \ROMsize\leq(\FOMsize-\ell)$, it holds that $\mathrm{det}(\bmK^{\top}(\Jcheck_1-\Rcheck_1)^{-1}\bmK)\neq 0$.
			Finally, we select $\bmK\in\matset_{(\Jcheck_1-\Rcheck_1)^{-1}}$, set $\bmU= \bmW\bmK$ and we can show that $[\bmB,\bmU]\in\matset_{(\bmJ-\bmR)^{-1}}$ by the following equation
			\begin{align*}
				&\mathrm{det}\left(
				\begin{bmatrix}
					\bmB,\bmW\bmK
				\end{bmatrix}^{\top}
				(\bmJ-\bmR)^{-1}
				\begin{bmatrix}
					\bmB,\bmW\bmK
				\end{bmatrix}
				\right)
				\\
				=
				&
				\mathrm{det}\left(
				\begin{bmatrix}
					\bmB^{\top}	(\bmJ-\bmR)^{-1}\bmB&
					\bmB^{\top}	(\bmJ-\bmR)^{-1}\bmW\bmK\\
					\bmK^{\top}\bmW^{\top}(\bmJ-\bmR)^{-1}\bmB&
					\bmK^{\top}(\bmW^{\top}(\bmJ-\bmR)^{-1}\bmW)\bmK
				\end{bmatrix}
				\right)
				\\
				=
				&
				\mathrm{det}\left(
				\begin{bmatrix}
					\bmB^{\top}	(\bmJ-\bmR)^{-1}\bmB&
					\bmB^{\top}	(\bmJ-\bmR)^{-1}\bmW\bmK\\
					\bzero&
					\bmK^{\top}(\bmJ_1-\bmR_1)\bmK
				\end{bmatrix}
				\right)
				\\
				=
				&
				\mathrm{det}\left(\bmB^{\top}	(\bmJ-\bmR)^{-1}\bmB\right)
				\mathrm{det}\left(\bmK^{\top}(\Jcheck_1-\Rcheck_1)^{-1}\bmK\right)
				\neq
				0.
\end{align*}							
\end{proof}					

\end{document}
